% Generated draft: requires mathematical and source review.
\documentclass[11pt]{article}
\usepackage[margin=25mm]{geometry}
\usepackage{fontspec}
\setmainfont{Latin Modern Roman}
\usepackage{amsmath,amssymb,mathtools,unicode-math}
\setmathfont{Latin Modern Math}
\usepackage{enumitem,needspace,xurl,hyperref}
\setlength{\parindent}{0pt}
\setlength{\parskip}{5pt}
\setlength{\emergencystretch}{3em}
\newcommand{\N}{\mathbb{N}}
\newcommand{\Z}{\mathbb{Z}}
\newcommand{\Q}{\mathbb{Q}}
\newcommand{\R}{\mathbb{R}}
\newcommand{\C}{\mathbb{C}}
\newcommand{\F}{\mathbb{F}}
\newcommand{\A}{\mathbb{A}}
\newcommand{\PP}{\mathbb{P}}
\newcommand{\E}{\mathbb{E}}
\newcommand{\eps}{\varepsilon}
\newcommand{\dd}{\,\mathrm{d}}
\begin{document}
\section*{UnsolvedMath ID: 2800101}
\textbf{Source problem code:} catalogue entry
\subsection{Definitions and mathematical statement}
Let \(n\ge1\), let \(\Sigma\) be a positive-semidefinite real symmetric \(n\times n\) matrix, and let \(X\sim N(0,\Sigma)\). For an orthonormal basis \(B=(b_1,\ldots,b_n)\), put
\[
N_B=\max_{1\le i\le n}|\langle X,b_i\rangle|,\qquad
\Phi_\Sigma(B)=\mathbb E N_B.
\]
A Karhunen--Lo\`eve (KL) basis \(B_{\mathrm{KL}}\) is an orthonormal eigenbasis of \(\Sigma\). If its eigenvalues are \(\lambda_i\), then its coordinates have the law \((\sqrt{\lambda_i}Z_i)_{i=1}^n\), with independent standard normal \(Z_i\).

The stored statement asks whether \(B_{\mathrm{KL}}\) minimizes \(\Phi_\Sigma\), namely whether
\[
\Phi_\Sigma(B_{\mathrm{KL}})\le\Phi_\Sigma(B)\quad\text{for every }B.
\]
The supplied research summary explicitly says the dataset has this extremum reversed and that the intended Mallat--Zeitouni conjecture is instead
\[
\forall n,\Sigma,B,\qquad
\Phi_\Sigma(B)\le\Phi_\Sigma(B_{\mathrm{KL}}),
\]
i.e. KL maximizes the expected largest absolute coordinate. Both formulations are displayed to preserve the material conflict; the maximum version is a supplied correction, not a silent replacement of the registry statement. Degenerate covariance matrices are not excluded.

The summary says the minimum version is already refuted in dimension \(2\), claims an exact dimension-\(2\) proof for the maximum version and a variance-equalizing minimizer, and refers to a ``closed form above'' that is absent from the record. Its numerical verification claim does not supply a proof artifact. It reports the exact maximum conjecture open in general, constant-factor approximate optimality attributed to Litvak--Tikhomirov (2018), and a \(1+O(1/\sqrt d)\) factor attributed to Liu (2026). The parameter \(d\), asymptotic regime, and precise hypotheses of that last estimate are not specified. These are unverified supplied progress claims.
\subsection{Short English statement}
Does an eigenbasis of a Gaussian covariance matrix maximize the expected largest coordinate magnitude? The stored question says minimize, while its research summary explicitly corrects that direction and reports partial progress.
\subsection{Sources}
\begin{itemize}
\item Local UnsolvedMath registry, permanent ID \(2800101\), exact source code \texttt{catalogue entry}, ``10 Lectures and 42 Open Problems --- Mallat-Zeitouni Gaussian-basis problem.'' Proposed-by metadata: Afonso S. Bandeira, 2015.
\item Bandeira, \emph{42 Open Problems in Data Science} (2016), Open Problem 1.1; extracted author-blog source \url{https://afonsobandeira.wordpress.com/2014/07/02/an-interesting-problem-by-mallat-and-zeitouni/}. HTML extraction; recorded access 2026-07-29. The URL dates the blog entry to 2014; list year and proposed-year metadata are retained separately.
\item Supplied research summary names Litvak--Tikhomirov (2018) and Liu (2026), without complete citations. Its claimed new dimension-\(2\) argument and referenced closed form are not included.
\end{itemize}
\end{document}
